\documentclass[11pt,letterpaper]{article}

\usepackage[T1]{fontenc}
\usepackage[utf8]{inputenc}
\usepackage[spanish,es-nodecimaldot]{babel}
\usepackage[margin=2.5cm]{geometry}
\usepackage{ragged2e}

\newcommand{\citydate}{Monterrey, Nuevo Leon, a 17 de junio de 2026}
\newcommand{\recipient}{Comite de Admision}
\newcommand{\programname}{Maestria en Gestion Administrativa}
\newcommand{\institutionname}{Instituto Tecnologico Superior de Cajeme (ITESCA)}
\newcommand{\subjectline}{Solicitud de ingreso a la Maestria en Gestion Administrativa}
\newcommand{\applicantname}{Martin Jonathan De la Cruz Munoz}

\setlength{\parindent}{0pt}
\setlength{\parskip}{0.55em}
\linespread{1.05}\selectfont

\begin{document}
\begin{flushright}
\citydate
\end{flushright}

\vspace{1.2em}
{\bfseries \recipient\par}
Presente.

\vspace{0.8em}
{\bfseries Asunto:} \subjectline

\vspace{0.8em}
\justifying
Por medio de la presente, expongo mi interes en ingresar a la \textbf{\programname} del \textbf{\institutionname}. Mi motivacion principal es fortalecer mi preparacion para analizar, disenar y conducir soluciones administrativas con enfoque estrategico, tecnico y humano, especialmente en organizaciones que requieren mejorar sus procesos, su productividad y su capacidad de adaptacion.

Mi formacion como ingeniero industrial y mi experiencia en areas vinculadas con procesos, automatizacion, docencia y solucion de problemas me han permitido observar que la mejora tecnica de una organizacion solo alcanza resultados sostenibles cuando se integra con una gestion administrativa solida. En la practica, los retos de calidad, competitividad, control de recursos y toma de decisiones exigen una vision integral que conecte la operacion con la planeacion y el uso inteligente de la informacion.

La MGA representa para mi una oportunidad pertinente para consolidar competencias en gestion organizacional, administracion, innovacion, calidad y competitividad. Mi interes es adquirir herramientas metodologicas que me permitan diagnosticar problemas administrativos, estructurar proyectos de mejora, evaluar alternativas con criterios objetivos y proponer intervenciones viables para empresas e instituciones que buscan fortalecer su desempeno.

Uno de mis objetivos academicos y profesionales es desarrollar una linea de trabajo orientada a la mejora de procesos administrativos y productivos mediante el uso de datos, indicadores de desempeno y metodologias de mejora continua. Considero que este enfoque puede aportar valor a organizaciones industriales, de servicios y educativas que requieren ordenar su operacion y tomar decisiones con mayor fundamento.

Cuento con disposicion para el aprendizaje autonomo, disciplina de estudio, capacidad de analisis y experiencia en la comunicacion de conocimientos. Asumo el compromiso de participar con responsabilidad en las actividades academicas del programa y aportar una perspectiva practica, etica y orientada a resultados. Agradezco la atencion brindada a mi solicitud y reitero mi interes en formar parte de la \programname\ de ITESCA.

\vspace{1.8em}
Atentamente,

\vspace{3.0em}
\rule{0.58\textwidth}{0.4pt}\\[-0.2em]
\textbf{\applicantname}\\
Aspirante

\end{document}